% Margenes
\geometry{
  a4paper,
  inner=2.5cm,
  outer=2cm,
  top=3.2cm,
  bottom=3cm,
  headheight=2.0cm,
  headsep=0.35 cm,
  footskip=1.5cm
}

% Espaciado general del texto
% \spacing{1.5}
\setstretch{1.5}

% Eliminar espaciado de enumerate e itemize
\setlist[enumerate,itemize]{itemsep=0pt, topsep=0pt}

% Estilo para páginas de título (como \maketitle)
\fancypagestyle{plain}{
    \fancyhf{}                          % Limpia pie y encabezado de paginas
    \renewcommand{\headrulewidth}{0pt}  % Quita la línea de arriba
    
    % Nombre abajo: En impares a la Izquierda (LO), en pares a la Derecha (RE)
    % Nombre SIEMPRE a la Izquierda (L)
    
    % Número SIEMPRE a la Derecha (R)
    \fancyfoot[R]{\textbf{Pág. N° \thepage}}
}

% Estilo general para el documento
\pagestyle{fancy}
    \fancyhf{}
    \renewcommand{\headrulewidth}{0pt}
    
    % --- ENCABEZADOS (LOGOS) ---
    % Logo UNaM a la Izquierda (Left) en páginas pares (E) e impares (O)
    \fancyhead[LE,LO]{\includegraphics[height=1.7 cm]{img/logo_UNaM.png}} 
    
    % Logo FI a la Derecha (Right) en páginas pares (E) e impares (O)
    \fancyhead[RE,RO]{\includegraphics[height=1.6 cm]{img/logo_FI_grande.png}}
    
    % --- PIE DE PÁGINA (TUS DATOS) ---
    % Nombre en la esquina interna del libro
    % Nombre SIEMPRE a la Izquierda (L)
    %\fancyfoot[L]{\textbf{Wlodeck, Fabricio Joaquín}}
    % Número SIEMPRE a la Derecha (R)
    \fancyfoot[R]{\textbf{Pág. N° \thepage}}

% Texto parte de las longtable (tablas de varias páginas)
\DefTblrTemplate{contfoot-text}{default}{Continua en la siguiente página $\cdots$}
\DefTblrTemplate{conthead-text}{default}{(Continuación)}

% Definición del espaciado de indice
\setlength{\cftbeforetoctitleskip}{-10 pt}       % Espacio antes de indice general
\setlength{\cftaftertoctitleskip}{15 pt}         % Espacio despues de indice general

% Cambiar el tamaño de la palabra "Índice"
\renewcommand{\cfttoctitlefont}{\Huge\bfseries} % Tamaño Huge y negrita

% Usar el punto como separador decimal
\decimalpoint

\begin{comment}
    \titleformat{nombre de la parte a modificar}
        {Formato de la sección}    
        {Número y separador}                             
        {Separación entre numero y tituli}                  
        {Código adicional}
        [Otra configuración, por ejemplo línea]
\end{comment}

% Redefinición del formato de \chapter
\titleformat{\chapter}[hang]
    {\normalfont\huge\bfseries} % sin \centering → alineado a la izquierda
    {\thechapter.}              % número con punto
    {0.25 cm}
    {\MakeUppercase}
    \titlespacing*{\chapter}{0 pt}{-45 pt}{15 pt} % {izq}{antes}{después}

% Redefinición del formato de \section
\titleformat{\section}[block]
    {\normalfont\Large\bfseries}    
    {}                    
    {0 cm}                         
    {\MakeUppercase}                
    [\vspace{0.5ex}\titlerule]       

% Redefinición del formato de \subsection
\titleformat{\subsection}
    {\normalfont\large\bfseries}    % Formato de la subsección
    {}                              % Elimina el número
    {0 pt}                          % Sangría
    {}                              % Código adicional (opcional)

% Redefinición del formato de \subsubsection
\titleformat{\subsubsection}
    {\normalfont\normalsize\bfseries}   % Formato de la subsubsección
    {}                                  % Elimina el número
    {0 pt}                              % Sangría
    {}                                  % Código adicional (opcional)

% Elimina la numeración del indice
    
    % Secciones
    \renewcommand{\cftsecpresnum}{}         % Elimina el prefijo de numeración
    \renewcommand{\cftsecaftersnum}{}       % Elimina el sufijo de numeración
    \renewcommand{\cftsecnumwidth}{0 pt}    % Ancho de la numeración (0pt para eliminarla)
    \renewcommand{\thesection}{}            % Elimina el número de sección en el índice
    
    % Subsecciones
    \renewcommand{\cftsubsecpresnum}{}          % Elimina el prefijo de numeración
    \renewcommand{\cftsubsecaftersnum}{}        % Elimina el sufijo de numeración
    \renewcommand{\cftsubsecnumwidth}{0 pt}     % Ancho de la numeración (0pt para eliminarla)
    \renewcommand{\thesubsection}{}             % Elimina el número de subsección en el índice
    
    % Subsubsecciones
    \renewcommand{\cftsubsubsecpresnum}{}       % Elimina el prefijo de numeración
    \renewcommand{\cftsubsubsecaftersnum}{}     % Elimina el sufijo de numeración
    \renewcommand{\cftsubsubsecnumwidth}{0 pt}  % Ancho de la numeración (0pt para eliminarla)
    \renewcommand{\thesubsubsection}{}          % Elimina el número de subsubsección en el índice